% !TeX TS-program = xelatex
%
% Copyright (C) 2003-2025 Javier Bezos http://www.texnia.com
%
% This file may be distributed and/or modified under the conditions of
% the MIT License. A version can be found at the end of this file.
%
% Repository: https://github.com/jbezos/enumitem
%

\documentclass[a4paper]{ltxguide}

\raggedright
\parskip=.8ex
\advance\oddsidemargin-.5cm
\advance\textwidth1cm
\addtolength{\textheight}{2cm}
\addtolength{\topmargin}{-1.5cm}

\usepackage{xcolor,bera}
\usepackage{xeCJK}
\setCJKmainfont{Source Han Serif SC}
\setCJKsansfont{Source Han Sans SC}
\setCJKmonofont{Source Han Sans SC}
\xeCJKsetup{CJKmath=true}

\definecolor{notes}{rgb}{.75, .3, .3}%

 \usepackage{hyperref}
 \hypersetup{hidelinks}
\makeatletter
\newenvironment{desc}
  {\if@nobreak
     \vskip-\lastskip
     \vspace*{-2.5ex}%
   \fi
   \decl}
  {\enddecl}
\def\@begintheorem#1#2{%
  \list{}{}%
  \global\advance\@listdepth\m@ne
  \item[{\sffamily\bfseries\color{notes}\MakeUppercase{#1}}]}%
\makeatother
\newtheorem{warning}{警告}
\newtheorem{note}{注意}
\newtheorem{example}{示例}
\makeatother

\usepackage{framed}
\definecolor{shadecolor}{rgb}{0.96,0.96,0.93}
\let\bblxv\verbatim
\let\bblexv\endverbatim
\def\verbatim{\begin{shaded*}\bblxv\vskip-\baselineskip\vskip2.5\parsep}
\def\endverbatim{\bblexv\vskip-2\baselineskip\end{shaded*}}

\newcommand\New[1]{%
  \colorbox[rgb]{.87, .9, .83}{新特性 #1}\enspace\ignorespaces}

\usepackage{hyperref}

\title{使用 \textsf{enumitem} 宏包\\定制列表}

\author{Javier Bezos\footnote{问题报告、评论和建议请发送至 \href{http://www.texnia.com/enumitem.html}%
{\texttt{http://www.texnia.com/enumitem.html}}。}}

\date{Version 3.11\\2025/02/06}

\IfFileExists{enumitem.sty}{\usepackage{enumitem}}{}
\IfFileExists{layouts.sty}{\usepackage{layouts}}{}

% Fix enumitem postlabel with CJK labels
\makeatletter
\let\enit@postlabel@i\@gobble
\makeatother{}

\addtolength{\topmargin}{-3pc}
\addtolength{\textwidth}{6pc}
\addtolength{\oddsidemargin}{-2pc}
\addtolength{\textheight}{7pc}

\raggedright
\parindent1.8em
\parskip0pt

\begin{document}

\maketitle
\tableofcontents

\newpage

\vspace*{1cm}

\begin{note}
  与 3.5 版相比的变更和新功能以 \New{X.X}\hspace{-.5em} 标出。最新功能可能仍不稳定。请将发现的问题报告至
  \texttt{https://github.com/jbezos/enumitem/issues}，这比在邮件列表或网络论坛上抱怨更有帮助。欢迎提交 fork 和 pull request。
\end{note}

\begin{warning}
  3.6 版引入了两个新键：|left| 和 |first|。如果你的文档中定义了使用这些名称的系列，将会报错。只需重命名它们，或设置宏包选项 |series=override|。
\end{warning}

\section{简介}

本宏包提供了你可能需要的绝大多数灵活性，用于定制三种基本列表环境（|enumerate|、|itemize| 和 |description|）以及设计自己的列表，采用 |<键>=<值>| 语法：
\begin{itemize}
\item 精美的标签和引用，包括与 \textsf{enumerate} 宏包类似的语法。

\item 标签定位的替代方式，引入了新参数（|labelindent|）和类似制表位的设置（|left|）。

\item 设置可全局应用，也可仅应用于三种列表类型之一，甚至单个列表（包括 |\topsep|）。

\item 行内列表。

\item 多种描述样式（同时修复了一些间距问题）。

\item 正确格式化的 |trivlist|。

\item 将多个列表组合为一个单元处理，以及计数器续编。

\end{itemize}

接口中使用了一种"继承"机制。你可以全局设置列表的行为，然后覆盖 |enumerate| 的若干参数，再在某个特定实例中覆盖少数参数。值将按层级逐级查找。

本宏包扩展了列表的语法，允许使用可选参数，其中提供了一组 |key=value| 形式的参数。这些键与众所周知的列表参数等价。请参阅 \LaTeX{} 手册了解它们的说明。下一节解释 \textsf{enumitem} 提供的扩展。
\begin{itemize}
\item
垂直间距：
\begin{itemize}
\setlength{\itemsep}{0pt}
\item |topsep|
\item |partopsep|
\item |parsep|
\item |itemsep|
\end{itemize}
\item
水平间距：
\begin{itemize}
\setlength{\itemsep}{0pt}
\item |leftmargin|
\item |rightmargin|
\item |listparindent|
\item |labelwidth|
\item |labelsep|
\item |itemindent|
\end{itemize}
\end{itemize}

\begin{example}
  一个简单的例子：
\begin{verbatim}
\begin{itemize}[itemsep=1ex, leftmargin=1cm]
\end{verbatim}
\end{example}

\begin{note}
  查看这些参数如何工作的一个方法是使用 |layouts| 宏包（手册名为 |layman.pdf|）。
  \ifx\listdiagram\notundefined\else
    参见图 \ref{f.drawlist}。
    \begin{figure}
      \centering
      \listdiagram
      \caption{列表参数}\label{f.drawlist}
    \end{figure}
  \fi
\end{note}

\section{快速参考}

一些常用设置。详见下文手册。

\begin{itemize}
\def\verbatim{\begin{shaded*}\bblxv}
\def\endverbatim{\bblexv\vskip-1.6\baselineskip\end{shaded*}}

\item 完全去除列表中的垂直间距：
\begin{verbatim}
\begin{enumerate}[nosep]
\end{verbatim}

\item 完全去除\textit{所有}列表中的垂直间距：
\begin{verbatim}
\setlist{nosep}
\end{verbatim}

\item 让标签从边距开始，条目文字从当前 parindent 开始：
\begin{verbatim}
\begin{enumerate}[left=0pt .. \parindent]
\end{verbatim}

\item 像 \textsf{enumerate} 那样配置标签：加载宏包选项 |shortlabels|，然后作为第一个元素写入你的
label:
\begin{verbatim}
\begin{enumerate}[(1)]
\end{verbatim}

\item 在"暂停"后继续上一个列表：
\begin{verbatim}
\begin{enumerate}[resume*]
\end{verbatim}

\item 行内使用三个基本列表：只需添加宏包选项 |inline|，然后使用环境 |enumerate*|、|itemize*| 和 |description*|。

\item 设置带括号的数字标签，但交叉引用不带括号：
\begin{verbatim}
\begin{enumerate}[label=(\arabic*), ref=\arabic*]
\end{verbatim}

\end{itemize}

\section{键}

本节介绍显示列表中的键。其中大多数也可用于行内列表，后者还有更多可用键（见 \ref{s.inline}）。

\begin{warning}
  如果值完全被花括号包围，它们会被剥离。如果你需要花括号，必须写成双份。这是 \textsf{keyval} 的默认行为，\textsf{enumitem} 只是模拟它。
\end{warning}

\subsection{标签与交叉引用格式}

\begin{desc}
|label=<commands>|
\end{desc}

设置当前层级使用的标签。一组不带参数的 |\alph|、|\Alph|、|\arabic|、|\roman| 和 |\Roman| 的星号版本代表 |enumerate| 中的当前计数器。\footnote{实际上，星号目前就是参数，但未来可能改变。请将它们视为星号变体并遵循相应语法。} 它也适用于 |\value|（前提是不需要计算最宽标签，或使用了 |widest*|，见下文）。

\begin{note}
  如果你想像 \textsf{enumerate} 宏包那样设置标签，请使用"短标签"（见第 \ref{s.short} 节）。
\end{note}

\begin{example}
  以下代码输出 \textit{a})、\textit{b}) 等（这是西班牙的标准样式，以前芝加哥也使用）：
\begin{verbatim}
\begin{enumerate}[label=\emph{\alph*})]
\end{verbatim}
\end{example}

\begin{warning}
  |label| 的值是移动参数，脆弱命令必须保护，\textit{但}计数器除外。因此，|\value| 的使用有些棘手，因为 |\the| 或 |\ifnum| 期望一个实际值，而在处理 |label| 以将内部的 |*| 替换为带计数器参数的形式时并非如此。最佳解决方案通常是借助 |\AddEnumerateCounter| 将逻辑封装在一个新的"计数器"中。\footnote{诚然这有些复杂。更好的方法正在开发中。}
\end{warning}

\begin{example}
  一个更花哨的例子（看起来不美观，但仅用于展示可能性；需要 \textsf{color} 和 \textsf{pifont}）：
\begin{verbatim}
\begin{enumerate}[label=\protect\fcolorbox{blue}{yellow}{\protect\ding{\value*}}]
\end{verbatim}
\end{example}

\begin{desc}
|label*=<commands>|
\end{desc}

类似 |label|，但其值追加到父标签之后。例如，以下定义了一个 |legal| 列表（1.、1.1.、1.1.1. 等）：
\begin{verbatim}
\newlist{legal}{enumerate}{10}
\setlist[legal]{label*=\arabic*.}
\end{verbatim}

\begin{desc}
|ref=<commands>|
\end{desc}

默认情况下，|label| 同时设置交叉引用和 |\the...| 的形式（覆盖父层级中的设置），但你可以用此键定义不同的格式。例如，去掉右括号：
\begin{verbatim}
\begin{enumerate}[label=\emph{\alph*}), ref=\emph{\alph*}]
\end{verbatim}

\begin{note}
  在 |label| 和 |ref| 中，计数器都可以照常使用。因此，如果当前层级是第二级：
\begin{verbatim}
\begin{enumerate}[label=\theenumi.\arabic*.]
\end{verbatim}
  或
\begin{verbatim}
\begin{enumerate}[label=\arabic{enumi}.\arabic*.]
\end{verbatim}  
\end{note}

\begin{note}
  |label| 不会\textit{累积}形成引用。如果你想要，比如第一级为 1)、第二级为 \textit{a})，从而得到类似 1.\textit{a} 的结果，必须用 |ref| 设置。你可以使用 |\ref{level1}.\ref{level2}| 配合适当的 |ref| 设置，但正如 Robin Fairbairns 在 \TeX{} FAQ 中指出的：
  \begin{quote}
  ……这样做既繁琐又容易出错。更重要的是，这并不可取，因为你构造了一个不灵活的视觉表示（你无法一次性改变列表元素的所有引用）。
  \end{quote}
  这很有道理，我建议遵循这个建议，但有时你可能想要这样的效果：
\begin{verbatim}
... subitem \ref{level2} of item \ref{level1} ...
\end{verbatim}
\end{note}

\begin{warning}
  |ref| 的值是移动参数，脆弱命令必须保护，\textit{但}计数器除外。
\end{warning}

\begin{desc}
|font=<commands>|\qquad|format=<commands>|
\end{desc}

设置标签字体。当通过 |\item| 的可选参数更改标签时，以及在 \texttt{description} 中很有用。|<commands>| 中最后一个命令可以接受条目标签作为参数。在 \texttt{description} 类设置生效时，你可能需要以 |\normalfont| 开头。同义词是 \texttt{format}。实际上，此键可用于在每个 |\item| 处、标签之前执行的任何内容。

\begin{desc}
|align=left|\qquad |align=right|\qquad |align=parleft|
\end{desc}

标签如何对齐（相对于标签框边缘）。有三个可选值：|left|、默认的 |right| 和 |parleft|（一个宽度为 |\labelwidth| 的 parbox，左对齐文本）。控制标签间距的参数应正确设置，手动设置或更方便地使用 |*| 设置（见下文）：
\begin{verbatim}
\begin{enumerate}[label=\Roman*., align=left, leftmargin=*]
\end{verbatim}
当标签框应为其自然宽度时，使用 |left|。

\begin{desc}
|\SetLabelAlign{<value>}{<commands>}|
\end{desc}

可以用 |\SetLabelAlign| 定义新的对齐类型（或重新定义现有类型）；预定义值等价于：
\begin{verbatim}
\SetLabelAlign{right}{\hss\llap{#1}}
\SetLabelAlign{left}{#1\hfil}
\SetLabelAlign{parleft}{\strut\smash{\parbox[t]\labelwidth{\raggedright##1}}}
\end{verbatim}

\begin{example}
  虽然主要用于对齐，但此命令还有其他用途（例如提供的 |parleft|）。例如，以下代码将所有 |align=right| 的标签设为上标：
\begin{verbatim}
\SetLabelAlign{right}{\hss\llap{\textsuperscript{#1}}}
\end{verbatim}
  当然，也可以使用新名称。
\end{example}

\begin{note}
  如果定义中的最后一项是 skip（通常是 |\hfil|），有时会被 |description| 移除。如果出于某种原因你想避免这种情况，只需在末尾加 |\null|。
\end{note}

\begin{note}
  如果你想让 \texttt{align} 和 \texttt{font} 的内部设置被忽略，可以在 \texttt{before} 中覆盖 \textsf{enumitem} 对 |\makelabel| 的定义：
\begin{verbatim}
\begin{description}[before={\renewcommand\makelabel[1]{\ref{##1}}}]
\end{verbatim}
  或者，定义一个宏并用 |\let|。
\end{note}

\subsection{标签的水平间距}

当前层级左边距中的水平空间按以下方式分配：\footnote{诚然，这些图并不完全清晰，我打算在未来版本中改进。}
\begin{center}
\begin{tabular}{cc}
\fbox{\fbox{\strut \texttt{labelindent}}
  \fbox{\strut \texttt{labelwidth}}
  \fbox{\strut \texttt{labelsep} $-$ \texttt{itemindent}}}
&
\fbox{\strut\texttt{itemindent}}\\
\texttt{leftmargin}
\end{tabular}
\end{center}
这里 |labelindent| 是 \textsf{enumitem} 引入的新参数，见下文。其余是标准 \LaTeX 中的参数。

实际上，布局更复杂，因为标签框（即 |labelwidth|）可能伸入边距，这意味着 |labelindent| 取负值。

\begin{note}
  由于 |\parindent| 在列表内部不被直接使用，而是内部设为 |\itemindent| 或 |\listparindent|，当作为参数值使用时，\textsf{enumitem} 返回全局值，即它在最外层列表之外的值。
\end{note}

\begin{note}
\New{3.6} 如果你觉得这些参数令人困惑，你不是一个人。你可以在第一个条目前（或在 |first| 中）写 |\DrawEnumitemLabel| 来可视化它们，它从上到下画4条线：|labelindent|、|labelwidth|、|labelsep|、|itemindent|（正细负粗）；|leftmargin| 用两条竖线标记。
\end{note}

\begin{desc}
|labelindent=<length>|\\
|\labelindent|
\end{desc}

此参数是 \textsf{enumitem} 新增的，表示从外层列表/文本边距到标签框左边缘的空白。这意味着存在冗余，因为其中一个参数依赖于其他参数，即它必须从其他值计算得出，如下所述。默认情况下，计算值是 |labelindent|，即使显式设置了某个值（默认为 0~pt）。因此，如果你将其设为某个值，很可能你想将某个其他参数设为 |!| 或 |*|，否则它会被忽略。

新增了一个计数器长度 |\labelindent|。

这五个参数的关系如下：
\[
|\leftmargin|+|\itemindent| = 
|\labelindent|+|\labelwidth|+|\labelsep|
\]

\begin{desc}
|left=<labelindent>|\\
|left=<labelindent> .. <leftmargin>|
\end{desc}

\New{3.6} 这是一个便捷键，用于快速设置标签最常见的布局。你可以将其视为一种带有两个制表位的"标尺"：标签起始位置和文本起始位置（两者都相对于正常边距）。仅给定 |<labelindent>| 时，左边距（即"文本起始位置"）用 |labelsep| 计算。它内部使用 |widest|，因此关于 |description| 的说明也适用于此：你可能需要更改计算出的参数（例如，配合 |align=left| 使用 |itemindent=*|）。

\begin{example}
  典型设置：
\begin{verbatim}
\begin{enumerate}[left= 0pt]
\begin{enumerate}[left= 0pt .. \parindent]
\begin{enumerate}[left= \parindent]
\begin{enumerate}[left= \parindent .. 2\parindent]
\begin{enumerate}[left= -\parindent .. 0pt]
\end{verbatim}
\end{example}

\begin{note}
  标签宽度设为默认的最宽值。如果列表中有 $\ge 10$ 的阿拉伯数字，你可能还需要设置 |widest|。
\end{note}

\begin{note}
|left=<labelindent>| 设置 |leftmargin=*|；
|left=<labelindent> .. <leftmargin>| 设置 |labelsep=*|。
\end{note}

\begin{desc}
|leftmargin=!|\qquad|itemindent=!|\qquad|labelsep=!|
\qquad|labelwidth=!|\qquad|labelindent=!|
\end{desc}

设置哪个值由其他值计算得出。默认是 |labelindent=!|，但注意有些键会设置另一个值（|wide| 和 description 的 |style|）。计算在\textit{所有}键都读完后进行。显式值不会丢失，因此对于以下层级设置：
\begin{verbatim}
leftmargin=2em
labelindent=1em,leftmargin=!
labelindent=!
\end{verbatim}
|leftmargin| 再次为 2em，而 |labelindent| 是计算参数。

\begin{note}
  在 |align=right|（默认）下，|labelindent=!| 和 |labelwidth=!| 在实际中表现相似。
\end{note}

\begin{desc}
|leftmargin=*|\qquad|itemindent=*|\qquad|labelsep=*|
\qquad|labelwidth=*|\qquad|labelindent=*|
\end{desc}

类似前面，但此外 |labelwidth| 首先设为当前标签的宽度，使用 |\arabic*| 中\textit{0}、|\roman*| 中\textit{viii}、|\alph*| 中\textit{m} 的默认值，大写形式类似（这些值可用 |widest| 更改，见下文）。示例：
\begin{verbatim}
\begin{itemize}[label=\textbullet, leftmargin=*]
\begin{enumerate}[label=\roman*), leftmargin=*, widest=iii]
\begin{itemize}[label      = \textbullet,
                leftmargin = 2pc,
                labelsep   = *]
\begin{enumerate}[label       = \arabic*.,
                  labelindent = \parindent,
                  leftmargin  = 2\parindent, 
                  labelsep    = *]
\end{verbatim}

最有用的是 |labelsep=*| 和 |leftmargin=*|。前者使条目正文从固定位置开始（即 |leftmargin|），后者从取决于标签的可变位置开始（当然在一个列表内总是相同的）。

\begin{note}
  遗憾的是，\LaTeX{} 没有定义一个应用于所有列表的默认 |labelsep|——只是使用当前值。使用 \textsf{enumitem} 你可以为每个列表设置默认值，如下所述，因此如果你想确保 |labelsep| 在你的控制之下，只需：
\begin{verbatim}
\setlist{labelsep=.5em}
\end{verbatim}
\end{note}

\begin{note}
  |labelwidth=*| 和 |labelwidth=!| 是同义词。请谨慎使用，因为它们可能取负值，这没有意义（会显示警告）。
\end{note}

\begin{desc}
|widest=<string>|\qquad|widest*=<integer>|\qquad|widest|
\end{desc}

如需与 \texttt{*}-值配合使用。它覆盖最宽打印计数器的默认值。有时，如果列表不太长，|\alph| 的值 |a| 比默认的 |m| 更合理：
\begin{verbatim}
\begin{enumerate}[leftmargin=*, widest=a] % Assume standard 2nd level
\end{verbatim}
不带值时恢复默认。使用 |widest*| 时，字符串使用 |<integer>| 作为计数器值构建（例如，对于 |\roman|，|widest=viii| 和 |widest*=8| 相同）。

由于 |\value| 返回的不是字符串而是数字，|widest| 和 |*| 值不能与它一起使用。但 |widest*| 是数字，所以允许。

\New{3.6} 它也可用于 |itemize| 和 |description|。但由于后者在标签内部做了一些复杂的格式化，你可能需要更改计算出的参数（例如，配合 |align=left| 使用 |itemindent=*|）。

\begin{desc}
|labelsep*=<length>|
\end{desc}

记住 |labelsep| 跨越 |leftmargin| 的一部分和 |itemindent|（如果后者非零）。这常常令人困惑，因此提供了一个新键——使用 \texttt{labelsep*} 时值从左边距计算（它只是设置 |\labelsep| 然后加上 |\itemindent|，但此外后续对 |itemindent| 的更改也会被考虑）：
\begin{center}
\begin{tabular}{cc}
\fbox{\fbox{\strut \texttt{labelindent}}
  \fbox{\strut \texttt{labelwidth}}
  \fbox{\strut \texttt{labelsep*}}}
&
\fbox{\strut\texttt{itemindent}}\\
\texttt{leftmargin}
\end{tabular}
\end{center}

\begin{desc}
|labelindent*=<length>|
\end{desc}

类似 |labelindent|，但从当前列表的左边距计算，而不是从外层列表/文本的左边距计算。

\begin{example}
  第一种模式使标签与周围的 |\parindent| 对齐，同时条目正文根据标签和固定的 |labelsep| 缩进：
\begin{verbatim}
labelindent = \parindent,
leftmargin  = *
\end{verbatim}
  一个常见变体是使标签与周围文本对齐（记住如果 |labelindent| 不是计算参数，默认为 |0pt|）：
\begin{verbatim}
leftmargin = *
\end{verbatim}
  前者在第一层看起来更好，后者在后续层级更可取。可以轻松设置：
\begin{verbatim}
\setlist{leftmargin=*}
\setlist[1]{labelindent=\parindent} % Only the level 1
\end{verbatim}
\end{example}

\begin{example}
  第二种模式使条目正文与周围的 |\parindent| 对齐。在这种情况下（记住如果未设置，|labelindent| 是计算参数）：
\begin{verbatim}
leftmargin = \parindent
\end{verbatim}
\end{example}

\begin{example}
  第三种模式是标签与 |\parindent| 对齐，条目正文与 |2\parindent| 对齐：
\begin{verbatim}
labelindent = \parindent, 
leftmargin  = 2\parindent, 
itemsep     = *
\end{verbatim}
  同样，一个变体是标签与周围文本对齐，条目正文与 |\parindent| 对齐：
\begin{verbatim}
leftmargin = \parindent,
itemsep    = *
\end{verbatim}
\end{example}

\subsection{编号、停止与续编}

\begin{desc}
|start=<integer>|
\end{desc}
设置第一个条目的编号。

\begin{desc}
|resume|
\end{desc}

计数器从之前的 |enumerate| 继续，
instead of being reset to 1.
\begin{verbatim}
\begin{enumerate}
\item First item.
\item Second item.
\end{enumerate}
Text.
\begin{enumerate}[resume]
\item Third item 
\end{enumerate}
\end{verbatim}

这是局部完成的。如果你想要全局续编，参见下一节关于
series.

\begin{desc}
|resume*|
\end{desc}

类似 |resume|，但同时使用之前列表的选项。此选项必须限制在环境的可选参数中（这是它唯一有意义的地方）。应谨慎使用——如果你经常使用它，很可能你需要定义一个新列表（见 \ref{s.clone}）。允许其他键，在这种情况下保存的选项被当前列表中的选项覆盖（即 \texttt{resume*} 的位置无关紧要）。如果某个列表系列使用 \texttt{resume*}，选项取自第一个之前的列表，\texttt{start} 除外。

\begin{example}
  例如：
\begin{verbatim}
\begin{enumerate}[resume*, start=1] % 或 [start=1, resume*]
\end{verbatim}
  使用前一个 \texttt{enumerate} 中的键，但重启计数器。
\end{example}

\subsection{系列}

\begin{desc}
|series=<series-name>|\\
|<series-name>|\qquad|resume*=<series-name>|
\qquad|resume=<series-name>|
\end{desc}

继续列表的另一种方法是使用键 \texttt{series}，使它们表现为一个单元。带键 \texttt{series} 的列表被视为起始列表，其设置\textit{全局}存储，以便后续通过 \texttt{resume}/\texttt{resume*} 使用。所有这些键都接受系列名作为值（必须不同于现有键）：
\begin{itemize}
\item |resume=<series-name>| 仅继续系列中的条目编号，
\item |resume*=<series>| 还应用起始列表的设置，
\item |<series>|，即系列名作为键，是 |resume*=<series>| 的替代写法。
\end{itemize}

\begin{example}
  考虑：
\begin{verbatim}
\begin{enumerate}[label=\arabic*(a), leftmargin=1cm, series=l_after]
\item A
\item B
\end{enumerate}
\end{verbatim}
  得到：1(a)  2(a)。可以继续：
\begin{verbatim}
\begin{enumerate}[label=\arabic*(b), resume*=l_after]
                 % 或 [label=\arabic*(b), l_after]
\item A
\item B
\end{enumerate}
\end{verbatim}
  得到：3(b)  4(b)。（但如果你愿意，也可以用 |start=1|。）
\end{example}

注意你可以添加更多参数，它们在起始列表保存的参数之后执行，因此优先于后者——特别是，|resume*| 本身优先于起始列表中的 |start|（例如 |start=1|）。

\begin{note}
  每次启动一个系列时，内部会定义若干命令。因此，为避免浪费资源，不重叠的系列应使用相同名称。
\end{note}

\begin{warning}
  宏包未来可能引入新键，因此直接将 |<series>| 作为键是前向不兼容的潜在来源。然而，在键名中使用连字符或星号以外的非字母字符（例如 |!notes| 或 |m_steps|）以及大写字母和数字是安全的，因为 \textsf{enumitem} 永远不会使用它们。\New{3.7} 如果你定义了一些全小写名称的系列，而引入了新的冲突键，除了更改名称外，另一个选择是宏包选项 |series=override|（错误消息是 \texttt{Invalid series name `key'}）。使用它系列名称优先于预定义键——但仅在绝对必要时使用。
\end{warning}

\subsection{惩罚值}

\begin{desc}
|beginpenalty=<integer>|\qquad
|midpenalty=<integer>|\qquad |endpenalty=<integer>|
\end{desc}

分别设置列表开始处、条目之间和列表结束处的惩罚值。请参阅 \LaTeX{} 或 \TeX{} 手册了解惩罚值如何控制分页。与其他参数不同，列表开始时它们的值不会重置为默认值，因此会应用于子列表。

\subsection{注入代码}

\begin{desc}
|before=<code>| \qquad |before*=<code>|
\end{desc}

在列表开始前执行代码（更准确地说，在用于定义列表的 |list| 环境的第二个参数中）。非星号形式设置要执行的代码，覆盖任何先前值；星号形式将代码添加到现有代码中（在设置层级中，见下文，\textit{不}相对于外层列表/文本）。它可以包含例如规则和文本，但未经广泛测试。所有计算已完成，你可以访问和操作列表参数。

\begin{example}
  将两边距（左右）都设为最宽标签宽度：
\begin{verbatim}
\setlist{leftmargin=*, before=\setlength{\rightmargin}{\leftmargin}}
\end{verbatim}
\end{example}

\begin{desc}
|after=<code>|\qquad|after*=<code>|
\end{desc}

同上，但在列表即将结束前执行。

\begin{desc}
|first=<code>|\qquad|first*=<code>|
\end{desc}

\New{3.6} 同上，但作为列表正文中的第一件事，因此
\begin{verbatim}
\begin{itemize}[first=<code>]
\end{verbatim}
等价于
\begin{verbatim}
\begin{itemize}
<code>
\end{verbatim}

\begin{example}
  通过 |first| 你可以为显示元素定义自己的环境。一个简单的例子：
\begin{verbatim}
\newlist{letter}{itemize}{1}
\setlist[letter]{first=\item[]\itshape, rightmargin=\leftmargin}
\end{verbatim}
这里不需要 |label|，因为它不被使用。
\end{example}

\begin{desc}
|\EnumitemId|
\end{desc}

\New{3.7} 为帮助某些任务，每个列表被分配一个唯一数字标识符，由 |\EnumitemId| 返回。

\begin{example}
  以下示例展示如何将 |\label| 与 |\EnumitemId| 结合使用，并用 |after| 自动将列表标签宽度设为最宽者（前提是 ref 等价于标签）：\footnote{参见
  \texttt{https://tex.stackexchange.com/questions/29322/%
  how-to-make-enumerate-items-align-at-left-margin}.}
\begin{verbatim}
\SetEnumitemKey{widestlabel}
  {labelwidth = \widthof{\ref{enum-\EnumitemId}},
   after      = \label{enum-\EnumitemId}}
\end{verbatim}
然后只需使用键 |widestlabel|。
\end{example}

\begin{example}
  倒计数也是可行的，但稍微复杂一些，我们需要一些"外部"帮助。以下是一种可能的解决方案，但不是唯一的一种（很可能甚至不是最好的——例如，|start| 实际上是空操作）。
\begin{verbatim}
\newcounter{rev-enumi}   \newcounter{rev-enumii}
\newcounter{rev-enumiii} \newcounter{rev-enumiv}

\newcommand\revcounter[1]{%
  \setcounter{#1}{\value{#1}-2}%
  \stepcounter{rev-#1}%
  \gdef\afterrev{%
    \refstepcounter{rev-#1}%
    \setcounter{rev-#1}{0}}}

\AddEnumerateCounter\revcounter\revcounter{} % 第二个是占位符

\SetEnumitemKey{revRoman}
  {start  = \getrefnumber{enum-\EnumitemId}-1,
   label  = (\Roman*)\revcounter*,
   ref    = (\Roman*),
   after  = \afterrev
            \label{enum-\EnumitemId}}
\end{verbatim}
\end{example}

\subsection{描述样式}

|description| 中可用的一个键。
\begin{desc}
|style=<name>|
\end{desc}

设置描述\textit{样式}。|<name>| 可为以下之一：
\begin{description}
\item[|standard|] 类似标准类中的 |description|，尽管在其他类中可能略有不同。标签被boxed。设置 |itemindent=!|。

\item[|unboxed|] 很像标准 |description|，但标签不被boxed，以避免不均匀间距和过长标签无法换行。设置 |itemindent=!|。

\item [|nextline|] 如果标签不适合边距，文本在下一行继续，否则放在宽度为 |\leftmargin| $-$ |\labelsep| 的框中，即条目正文永远不会伸入左边距。设置 |labelwidth=!|。

\item[|sameline|] 类似 |nextline|，但如果标签不适合边距，文本在同一行继续。等同于 |style=unboxed,labelwidth=!|。

\item[|multiline|] 标签放在宽度为 |leftmargin| 的 parbox 中，必要时多行。等同于 |style=standard,align=parleft,labelwidth=!|。如果你修改它，记住 |align| 不能自由设置（出于内部优化）；特别是，水平框没有意义且可能引发错误。

\end{description}

\begin{warning}
\begin{enumerate}
\item 混合使用 boxed 和 unboxed 标签没有明确定义的行为。
\item 嵌套列表时所有组合都允许，但并非所有都有意义。
\item 不支持嵌套 |nextline| 列表（可以工作，但未来行为可能改变，因为当前行为并不符合预期）。
\end{enumerate}
\end{warning}

\subsection{紧凑列表}

\begin{desc}
|noitemsep|\qquad|nosep|
\end{desc}

键 |noitemsep| 消除条目和段落之间的间距（即 |itemsep=0pt| 和 |parsep=0pt|），而 |nosep| 消除所有垂直间距。\footnote{键 \texttt{nolistsep} 现已弃用，它引入了少量伸展，并非预期行为。}

\subsection{``Wide'' lists}

\begin{desc}
|wide|\\
|wide=<parindent>|
\end{desc}
使用这个便捷键，左边距为零，标签是文本的一部分——换句话说，条目看起来像普通段落。\footnote{\texttt{fullwidth} 已弃用。} 这里 |labelsep| 设置标签与第一个词之间的间距。等价于
equivalent to
\begin{verbatim}
align=left, leftmargin=0pt, labelindent=\parindent,
listparindent=\parindent, labelwidth=0pt, itemindent=!
\end{verbatim}
使用 |wide=<parindent>| 可以立即设置 |\parindent| 以外的值。当然，这些键也可以在 |wide| 之后覆盖；例如，记住左对齐标签如果比 |labelwidth| 宽会推动文本，你可以设置 |labelwidth=1.5em| 作为最小宽度，或者你可能更喜欢 |itemindent=*| 而不是 |itemindent=!|，后者将最小宽度设为最宽标签的宽度。在第2层你可能更喜欢 |labelindent=2\parindent|，依此类推。你可能还想将其与 |noitemsep| 或 |nosep| 组合使用。

\subsection{类似 \textsf{enumerate} 的标签}
\label{s.short}

\begin{desc}
|shortlabels|（宏包选项）
\end{desc}

使用宏包选项 \texttt{shortlabels} 你可以使用类似 \textsf{enumerate} 的语法，其中 |A|、|a|、|I|、|i| 和 |1| 分别代表 |\Alph*|、|\alph*|、|\Roman*|、|\roman*| 和 |\arabic*|。这主要是作为与 \textsf{enumerate} 宏包的一种兼容模式，因此适用以下特殊规则：如果（任意层级的）第一个选项不被识别为有效键，则被视为使用类似 \textsf{enumerate} 语法的标签。例如：
\begin{verbatim}
\begin{enumerate}[i), labelindent=\parindent, labelsep=*]
...
\end{enumerate}
\end{verbatim}
如果 |ref| 与 label 不同，你可能还需要设置它。

虽然可能不那么有用，在类似条件下你也可以在 itemize 环境中省略 |label=|：
\begin{verbatim}
\begin{itemize}[\textbullet]
...
\end{itemize}
\end{verbatim}

\begin{desc}
|\SetEnumerateShortLabel{<key>}{<replacement>}|
\end{desc}

使用此命令，你可以定义新键（或重新定义），这对于将 enumerate 适配到特定排版规则或扩展到非拉丁脚本特别有用。这里 |<replacement>| 包含计数器的星号版本之一。

\begin{example}
  例如：
\begin{verbatim}
\SetEnumerateShortLabel{i}{\textsc{\roman*}}
\end{verbatim}
  重新定义 |i|，使使用此键的条目以小型大写罗马数字编号。
\end{example}

\begin{note}
  键必须是单个字符。
\end{note}

\subsection{通用键与值}

\begin{desc}
|\SetEnumitemKey{<key>}{<replacement>}|
\end{desc}

使用此命令你可以创建自己的（无值）键。如此定义的键可以像其他键一样使用。

\begin{example}
  通过
\begin{verbatim}
\SetEnumitemKey{midsep}{topsep=3pt, partopsep=0pt}
\end{verbatim}
  你可以写
\begin{verbatim}
\begin{enumerate}[midsep]
\end{verbatim}
\end{example}

\begin{example}
  另一个例子是使用 \textsf{multicol} 的多列列表：
\begin{verbatim}
\SetEnumitemKey{twocol}{
  itemsep = 1\itemsep,
  parsep  = 1\parsep,
  before  = \raggedcolumns\begin{multicols}{2},
  after   = \end{multicols}}
\end{verbatim}
  这里，\texttt{itemsep} 和 \texttt{parsep} 的设置消除了伸展和收缩部分，在这种情况下改善了结果。当然，你可能需要定义一个新列表。
\end{example}

\begin{warning}
  本宏包未来可能引入新键，因此 |\SetEnumitemKey| 存在前向不兼容的风险。但在键名中使用连字符或星号以外的非字母字符（例如 |:name| 或 |2_col|）以及大写字母和数字是安全的，因为 \textsf{enumitem} 永远不会使用它们。
\end{warning}

\begin{desc}
|\SetEnumitemValue{<key>}{<string-value>}{<replacement>}|
\end{desc}
此命令为 |<key>=<value>| 对提供了进一步的抽象层。使用它你可以定义逻辑名称，它们被翻译为实际值。例如：
\begin{verbatim}
\SetEnumitemValue{label}{numeric}{\arabic*.}
\SetEnumitemValue{leftmargin}{standard}{\parindent}
\end{verbatim}
你可以这样写：
\begin{verbatim}
\begin{enumerate}[label=numeric, leftmargin=standard]
\end{verbatim}
因此，你可以将 |label=numeric| 的含义留到最终设计时决定。

\section{行内列表}
\label{s.inline}

行内列表是在段落中作为普通文本设置的"水平"列表。使用本宏包你可以创建行内列表，如下所述，通过 |\newlist|，它们有自己的标签和计数器。然而，通常使用标准列表的行内版本、采用相同标签模式就足够了——宏包选项 |inline| 就是做这个的。

\begin{warning}
  条目被boxed，因此浮动体丢失且不允许嵌套列表（记住许多显示元素被定义为列表）。链接可能指向错误位置，禁止使用行间公式，而且由于 \TeX{} 构建列表时做的优化，显式断词可能错误。\footnote{Knuth 式的"过早优化"？谁知道，但无论如何使用 Lua\TeX{} 断词应该是正确的。} 这个默认设置是有原因的，即此功能主要为短条目（几个词）设计，参数 |itemjoin*| 可用于逻辑标记。要克服这些限制，你可以设置 |mode=unboxed|，见下文。
\end{warning}

\begin{desc}
|inline| \qquad(package option)\\
\texttt{enumerate*}\qquad\texttt{itemize*}\qquad
\texttt{description*} \qquad(environments)
\end{desc}

使用宏包选项 \texttt{inline}，三个行内列表环境被定义：\texttt{enumerate*}、\texttt{itemize*} 和 \texttt{description*}。它们模拟 \textsf{paralist} 和 \textsf{shortlst} 的行为，即标签和设置与显示（即正常）列表 \texttt{enumerate}、\texttt{itemize} 和 \texttt{description} 分别共享（但记住恢复基于环境名称，而非列表类型）。这仅适用于用 \texttt{inline} 创建的那些——用 |\newlist| 创建的行内列表（如下所述）是独立的，使用自己的标签和设置。

\begin{note}
  注意如果你不需要标准列表的行内版本，则不需要 |inline|，而是用 |\newlist| 定义你自己的独立行内列表。
\end{note}

\begin{warning}
  由 |inline| 定义的这三个环境的设置与显示变体共享，因此不能用 |\newlist| 直接重定义。尝试这样做会引发晦涩错误。如果你需要独立设置，用 |\newlist| 定义而不用 |inline|。
\end{warning}

\begin{desc}
|itemjoin=<string>|\qquad|itemjoin*=<string>|
\qquad|afterlabel=<string>|
\end{desc}

格式通过键 \texttt{itemjoin}（默认是空格）和 \texttt{afterlabel}（默认是 |\nobreakspace|，即 |~|）设置。另一个键是 \texttt{itemjoin*}，如果设置，它在最后一个条目之前替代 \texttt{itemjoin}。

|itemjoin| 在垂直模式下被忽略（即在 unboxed 模式下，以及在引用、显示列表等之后）。

\begin{example}
  通过
\begin{verbatim}
before=\unskip{: }, itemjoin={{; }}, itemjoin*={{, and }}
\end{verbatim}
  条目之间使用以下标点：
  \begin{quote}
  Blah blah: (a) one; (b) two; (c) three, and (d) four. Blah blah
  \end{quote}
\end{example}

\begin{desc}
|mode=unboxed|\qquad|mode=boxed|
\end{desc}

如果在行内列表中使用浮动体、列表或行间公式很重要，请使用替代"模式"，通过 \texttt{mode=unboxed} 激活（默认是 \texttt{mode=boxed}）。使用它，浮动体可以自由使用，但错误放置的 |\item| 不会被捕获，且 \texttt{itemjoin*} 被忽略（日志中会写入关于此事实的警告）。

\section{全局设置}

也可以进行应用于所有这些列表的全局更改：
\begin{desc}
|\setlist[enumerate,<levels>]{<format>}|\\
|\setlist[itemize,<levels>]{<format>}|\\
|\setlist[description,<levels>]{<format>}|\\
|\setlist[<levels>]{<format>}|
\end{desc}
其中 |<level>| 是 |list| 中的列表层级（一个或多个），以及 |enumerate| 和 |itemize| 中的对应层级。\footnote{|\string\setenumerate|、|\string\setitemize| 和 |\string\setdescription| 已弃用。} 不带 |<levels>| 时，格式应用于所有层级。这里的 list 不是指任意列表，而是仅指本宏包处理的三个列表，以及被本宏包重定义或用 |\newlist| 定义的列表（见下文）。例如：
\begin{verbatim}
\setlist{noitemsep}
\setlist[1]{labelindent=\parindent} % << Usually a good idea
\setlist[itemize]{leftmargin=*}
\setlist[itemize,1]{label=$\triangleleft$}
\setlist[enumerate]{labelsep=*, leftmargin=1.5pc}
\setlist[enumerate,1]{label = \arabic*.,
                      ref   = \arabic*}
\setlist[enumerate,2]{label = \emph{\alph*}),
                      ref   = \theenumi.\emph{\alph*}}
\setlist[enumerate,3]{label = \roman*),
                      ref   = \theenumii.\roman*}
\setlist[description]{font=\sffamily\bfseries}
\end{verbatim}

这些设置按以下顺序读取：list、当前层级的 list、enumerate/itemize/description、以及当前层级的 enumerate/itemize/description；如果一个键多次出现且值不同，最后一个即最具体的值生效。如果我们用 \texttt{resume*} 恢复系列或列表，保存的键随后应用。最后，应用可选参数（\texttt{resume*} 除外）。

\LaTeX{} 提供了一组宏来更改这些参数中的许多，但用本宏包设置它们更一致，有时更灵活，代价是更"显式"（和冗长）。
列表规范可以包含变量和计数器，前提是它们可展开，且计数器支持 \textsf{calc}，因此如果你加载本宏包，可以写：
\begin{verbatim}
\newcount{toplist}
\setcount{toplist}{1}
\newcommand{\mylistname}{enumerate}
\setlist[\mylistname,\value{toplist}+1]{labelsep=\itemindent+2em]
\end{verbatim}
这允许借助循环定义列表。

\begin{warning}
  似乎没有办法捕获 |\setlist| 或 |\newlist| 中拼写错误的名称，会引发无意义的错误"缺少数字，按零处理"。
\end{warning}

\section{依赖字号的设置}
\label{s:sized}

\New{3.6} 对于依赖字号的设置，大多数情况下你只需要 |ex| 或 |em| 等相对单位。有时你可能想要离散的步骤，特殊语法允许这样做。

使用此功能需要以下宏包选项。

\begin{desc}
|sizes|（宏包选项）
\end{desc}

长度可以包含基于字号的设置，如下所示（第一个 \texttt{<} 之前的值是默认值）。
\begingroup
\makeatletter
\renewcommand\verbatim@font{\normalfont\ttfamily}
\begin{verbatim}
\setlist{
  topsep      = 20pt <-10> 6pt          <10-> 40pt,
  leftmargin  =      <-10> 0cm <10> 1cm <10-> 2cm ,
  rightmargin =      <-10> 0cm <10> 1cm <10-> 2cm ,
  }
\end{verbatim}
\endgroup

也接受名称：|script|、|tiny|、|footnote|、|small|、|normal|、|large|、|Large|、|LARGE|、|huge|、|Huge|（即必要时从 \LaTeX{} 名称中去掉 size）。例如：
\begingroup
\makeatletter
\renewcommand\verbatim@font{\normalfont\ttfamily}
\begin{verbatim}
\setlist{
  topsep      = 20pt <-normal> 6pt              <normal-> 40pt,
  leftmargin  =      <-normal> 0cm <normal> 1cm <normal-> 2cm ,
  rightmargin =      <-normal> 0cm <normal> 1cm <normal-> 2cm ,
  }
\end{verbatim}
\endgroup
单值优先于范围（即特定值优先于通用值）。在范围内，最后一个匹配生效。范围 |a-b| 表示 $$a \le \mbox{size} < b$$（包含下限，不包含上限）。这些规则在上面的示例中允许在逻辑位置设置 |10| 或 |normal|。记住标称字号并不总是实际字号——例如，使用选项 |11pt| 时，|\normalsize|（和 |normal|）是 10.95。你可以在一个值前加多个单限定符，如 \texttt{<}|small|\texttt{><}|normal|\texttt{>}|12pt|。\footnote{注意此语法紧密遵循 \texttt{\string\DeclareFontShape}，除了单值的优先级。} 第一个 \texttt{<..>} 之前的值被视为默认值。

\begin{note}
  出于效率原因，命名字号仅在加载 \textsf{enumitem} 时分配一次，假设它们由文档类设置，或先前加载的本地样式设置。
\end{note}

\begin{desc}
|\SetEnumitemSize{<name>}{<selector>}|
\end{desc}

\New{3.7} 如果在加载 \textsf{enumitem} 后修改了字号，或者你使用的文档类有非标准字号（甚至你只是想要其他名称），可以用以下工具设置或重置。

\begin{example}
一个简单例子：
\begin{verbatim}
\SetEnumitemSize{normal}{\normalsize}
\SetEnumitemSize{large}{\large}
\end{verbatim}
\end{example}

\begin{desc}
|\setlist|\texttt{\string<}|<size>|\texttt{\string>}%
   |[<names>,<levels>]{<keys/values>}|
\end{desc}

\New{3.7} 下面介绍的 |\setlist| 扩展，将字号优先级规则也应用于此（因此 |\setlist| 的顺序很重要），通过此过程定义的依赖字号的键优先于其他键。例如：
\begingroup
\makeatletter
\renewcommand\verbatim@font{\normalfont\ttfamily}
\begin{verbatim}
\setlist<-normal>[enumerate]{nosep}
\end{verbatim}
\endgroup

但每个 |\setlist| 中只接受一个字号限定符。

\section{克隆基本列表}
\label{s.clone}

\begin{desc}
|\newlist{<name>}{<type>}{<max-depth>}|\\
|\renewlist{<name>}{<type>}{<max-depth>}|
\end{desc}
三个列表可以被克隆，以便你定义表现类似它们的逻辑环境。要定义新列表（或重定义现有列表），使用 |\newlist|（或 |\renewlist|），其中 |<type>| 是 |enumerate|、|itemize| 或 |description|。行内列表的类型为 \texttt{enumerate*}、\texttt{itemize*} 和 \texttt{description*}。

\begin{note}
  行内 |<type>| 始终可用，即使没有宏包选项 |inline|，后者只是用这些名称定义对应类型的三个环境。
\end{note}
如果 |<type>| 是 |enumerate|，会定义一组名称为 |<name>i|、|<name>ii|、|<name>iii|、|<name>iv| 等的计数器（取决于 |<max-depth>|）。
然后你可以在标签中使用这些计数器；例如，如果你定义了名为 \texttt{steps} 的列表，可以这样定义标签：
\begin{verbatim}
label=\arabic{stepsii}.\arabic{stepsi}
\end{verbatim}

\begin{warning}
  不要为 |<max-depth>| 使用任意大的数字，以避免创建过多计数器和相关宏。
\end{warning}

\begin{warning}
  为与 \LaTeX 中的计数器命名方案保持一致，列表层级在内部也用罗马数字命名，即 |<list>i|、|<list>ii|、|<list>iii| 等。由于这个原因（计数器和列表名称都是），定义两个列表如 |books| 和 |booksi| 会导致意外结果（目前没有任何警告，应该修复）。
  % TODO - Perhaps I must change the internal names, so that 
  % this restriction doesn't apply to itemize and description.]
\end{warning}

\begin{desc}
|\setlist[<names>,<levels>]{<keys/values>}|\\
|\setlist*[<names>,<levels>]{<keys/values>}|
\end{desc}
创建列表后，你可以（实际上必须，至少设置标签）用 |\setlist| 设置新列表：
\begin{verbatim}
\newlist{ingredients}{itemize}{1}
\setlist[ingredients]{label=\textbullet}
\newlist{steps}{enumerate}{2}
\setlist[steps,1,2]{label=(\arabic*)}
\end{verbatim}
|\setlist| 可选参数中的名称指定设置应用于哪些列表，数字指定层级（支持 |calc|）。可以给出多个列表和/或多个层级，所有组合都会被设置；例如：
\begin{verbatim}
\setlist[enumerate,itemize,2,3]{...}
\end{verbatim}
不带数字（或0）表示所有层级，不带名称表示所有列表；不带可选参数表示所有列表的所有层级。
星号形式 |\setlist*| 将设置添加到先前的设置中。你可以将添加限制在特定字号，如第 \ref{s:sized} 节所述。必须注意后者优先级高于星号形式（即为某些字号添加的设置优先于为所有字号添加的设置，因此应用键的最具体值）。

\begin{desc}
|\setlistdepth{<integer>}|
\end{desc}

默认情况下，\LaTeX{} 限制5个嵌套层级，但克隆列表时这个值可能太短，因此你可能想设置新值。在第5层以下（或文档类定义的最深层级以下），使用最后一层的设置（即 |\@listvi|）。

\section{计数器详解}

\subsection{新计数器表示}

\begin{desc}
|\AddEnumerateCounter{<LaTeX command>}{<internal command>}{<widest label>}|
\end{desc}

"注册"一个计数器表示，使 \textsf{enumitem} 识别它。主要用于非拉丁文字，但在拉丁文字中也有用。

\begin{example}
以下示例定义了一个带有序数名称的新计数器：
\begin{verbatim}
\makeatletter
\def\ctext#1{\expandafter\@ctext\csname c@#1\endcsname}
\def\@ctext#1{\ifcase#1\or First\or Second\or Third\or
Fourth\or Fifth\or Sixth\fi}
\makeatother
\AddEnumerateCounter{\ctext}{\@ctext}{Second}
\end{verbatim}
\end{example}

\begin{note}
  计数器名称可以包含 |@|，即使它不是字母，也不会引发错误，如上例所示。
\end{note}
星号变体允许给出数字而不是字符串作为最宽标签。

\begin{example}
  如果最宽标签对应值2：
\begin{verbatim}
\AddEnumerateCounter*{\ctext}{\@ctmoreext}{2}
\end{verbatim}
\end{example}
如果表示不是普通字符串而是带样式的（例如小型大写），首选此变体。

\begin{example}
  俄语的例子：
\begin{verbatim}
\AddEnumerateCounter*{\asbuk}{\c@asbuk}{7}
\end{verbatim}
\end{example}

\subsection{重启 \texttt{enumerate}}

\begin{desc}
|\restartlist{<list-name>}|
\end{desc}

目前你可以通过以下方式获得文档中连续编号：
\begin{verbatim}
\setlist[enumerate]{resume}
\end{verbatim}
|\restartlist| 用于在文档中途重启计数器。例如，你可以在章节开始时发出 |\restartlist|，使每章都有连续编号。

它仅基于列表\textit{名称}，\textit{不}基于列表\textit{类型}，这意味着用宏包选项 \texttt{inline} 定义的 \texttt{enumerate*} 与 \texttt{enumerate} 不同，因为它的名称不同。

\section{宏包选项}

除了 |inline|、|ignoredisplayed|、|sizes|、|series=override| 和 |shortlabels|，还有以下选项。

\begin{desc}
|loadonly|
\end{desc}
使用此宏包选项，宏包被加载但三个列表不被重定义。你可以创建自己的列表，甚至重定义现有的。

\section{trivlist 问题}

\LaTeX{} 使用名为 |trivlist| 的 |list| 简化版本来设置显示元素，如 |center|、|tabbing|、|theorem| 等，即使概念上它们不是列表。不幸的是，|trivlist| 使用当前列表设置，这有一个奇怪的副作用：更改列表的垂直间距有时也会更改这些环境中的间距。
本宏包修改了 |trivlist|，使当前层级的默认设置（即相应 |clo| 文件设置的那些）被重新设置。在标准 \LaTeX{} 中这通常是多余的，但如果我们想微调列表，不重置默认值可能是个真正的问题（特别是如果你使用 |nosep| 选项）。
垂直间距的最小控制可通过以下方式实现\footnote{|\string\setdisplayed| 已弃用。}
\begin{itemize}
\item |\setlist[trivlist,<level>]{<keys/values>}|
\end{itemize}
但 |trivlist| 本身（不常直接使用）不接受可选参数。此功能并非旨在作为完整的 |trivlist| 格式化工具。
如果出于某种原因你不想更改 |trivlist| 并保留原始定义，可以使用宏包选项 |ignoredisplayed|。
\New{3.6} 另一方面，如果你还想将所有列表的更改应用于 trivlist，只需设置宏包选项 |includedisplayed|。

\section{示例}

\expandafter\ifx\csname setenumerate\endcsname\relax

请先安装宏包，然后重新排版本文档。

\else

在这些示例中，我们设置 |\setlist{noitemsep}|

\setlist{noitemsep}
\small

\newcommand{\newsample}{\vskip6pt\goodbreak\hrule height 1pt\vskip6pt}
\newcommand{\samplesep}{\vskip6pt\goodbreak\hrule\vskip6pt}
\newbox\vsep
\setbox\vsep\hbox{\vrule height 2ex depth 16ex width 1pt}
\dp\vsep0pt
\newcommand\showsep{\leavevmode\llap{\copy\vsep}}

\newsample

\begin{verbatim}
在拉曼恰的一个地方，其名我不愿提起，
不久前住着一位绅士
\begin{enumerate}[labelindent=\parindent,leftmargin=*]
  \item 一支长矛，
  \item 一面旧盾，
  \item 一匹瘦马，
  \item 一只猎犬。
\end{enumerate}
一锅肉多菜少的炖菜，多数晚上吃凉拌肉，
星期六吃煎羊肉……
\end{verbatim}

线显示 |labelindent|。 

\samplesep

\showsep 在拉曼恰的一个地方，其名我不愿提起，
不久前住着一位绅士
\begin{enumerate}[labelindent=\parindent,leftmargin=*]
\item 一支长矛,
\item 一面旧盾,
\item 一匹瘦马，
\item 一只猎犬.
\end{enumerate}
一锅肉多菜少的炖菜，多数晚上吃凉拌肉，
星期六吃煎羊肉……

\newsample

使用 |\begin{enumerate}[leftmargin=*] % labelindent 默认为 0pt|。 

线显示 |labelindent|。

\samplesep

\noindent\showsep\hskip\parindent 在拉曼恰的一个地方，其名我不愿提起，
不久前住着一位绅士
\begin{enumerate}[leftmargin=*]      
\item 一支长矛,
\item 一面旧盾,
\item 一匹瘦马，
\item 一只猎犬.
\end{enumerate}
一锅肉多菜少的炖菜，多数晚上吃凉拌肉，
星期六吃煎羊肉……

\newsample

使用 |\begin{enumerate}[leftmargin=\parindent]|。

线显示 |leftmargin|。

\samplesep

\showsep 在拉曼恰的一个地方，其名我不愿提起，
不久前住着一位绅士
\begin{enumerate}[leftmargin=\parindent]
\item 一支长矛,
\item 一面旧盾,
\item 一匹瘦马，
\item 一只猎犬.
\end{enumerate}
一锅肉多菜少的炖菜，多数晚上吃凉拌肉，
星期六吃煎羊肉……

\newsample

使用 |\begin{enumerate}[labelindent=\parindent,|\allowbreak
| leftmargin=*,|\allowbreak| label=\Roman*.,|\allowbreak
| widest=III,|\allowbreak| align=left]|.

线显示 |labelindent|。 Note 

\samplesep

\showsep 在拉曼恰的一个地方，其名我不愿提起，
不久前住着一位绅士
\begin{enumerate}[labelindent=\parindent, leftmargin=*,
                  label=\Roman*., widest=III, align=left]
\item 一支长矛,
\item 一面旧盾,
\item 一匹瘦马，
\item 一只猎犬.
\end{enumerate}
一锅肉多菜少的炖菜，多数晚上吃凉拌肉，
星期六吃煎羊肉……

\newsample

使用 |\begin{enumerate}[label=\fbox{\arabic*}]|。第一个条目为 \ref{i:first}

\samplesep

在拉曼恰的一个地方，其名我不愿提起，
不久前住着一位绅士
\begin{enumerate}[label=\fbox{\arabic*}]
\item \label{i:first}一支长矛,
\item 一面旧盾,
\item 一匹瘦马，
\item 一只猎犬.
\end{enumerate}
一锅肉多菜少的炖菜，多数晚上吃凉拌肉，
星期六吃煎羊肉……

\newsample

嵌套列表。

\samplesep

\begin{verbatim}
在拉曼恰的一个地方，其名我不愿提起，
不久前住着一位绅士
\begin{enumerate}[label=(\alph*), labelindent=\parindent,
     leftmargin=*, start=12]
\item 一支长矛，
\begin{enumerate}[label=(\alph{enumi}.\roman*), leftmargin=*, start=7]
\item 一面旧盾，
\end{enumerate}
\item 一匹瘦马，
\begin{enumerate}[label=(\alph{enumi}.\roman*), leftmargin=*, resume]
\item 一只猎犬。
\end{enumerate}
\end{enumerate}
一锅肉多菜少的炖菜，多数晚上吃凉拌肉，
星期六吃煎羊肉……
\end{verbatim}

在拉曼恰的一个地方，其名我不愿提起，
不久前住着一位绅士
\begin{enumerate}[label=(\alph*), labelindent=\parindent,
     leftmargin=*, start=12]
\item 一支长矛,
\begin{enumerate}[label=(\alph{enumi}.\roman*), leftmargin=*, start=7]
\item 一面旧盾,
\end{enumerate}
\item 一匹瘦马，
\begin{enumerate}[label=(\alph{enumi}.\roman*), leftmargin=*, resume]
\item 一只猎犬.
\end{enumerate}
\end{enumerate}
一锅肉多菜少的炖菜，多数晚上吃凉拌肉，
星期六吃煎羊肉……

\newsample

\begin{verbatim}
在拉曼恰的一个地方，其名我不愿提起，
不久前住着一位绅士
\begin{description}[font=\sffamily\bfseries, leftmargin=3cm,
    style=nextline]
  \item[\textnormal{堂吉诃德最先有的}] 一支长矛，
  \item[\textnormal{其次}] 一面旧盾，
  \item[\textnormal{再次}] 一匹瘦马，
  \item[\textnormal{最后，第四样}] 一只猎犬。
\end{description}
一锅肉多菜少的炖菜，多数晚上吃凉拌肉，
星期六吃煎羊肉……
\end{verbatim}

\samplesep

在拉曼恰的一个地方，其名我不愿提起，
不久前住着一位绅士
\begin{description}[font=\sffamily\bfseries, leftmargin=3cm,
    style=nextline]
\item[\textnormal{堂吉诃德最先有的}] 一支长矛,
\item[\textnormal{其次}] 一面旧盾,
\item[\textnormal{再次}] 一匹瘦马，
\item[\textnormal{最后，第四样}] 一只猎犬.
\end{description}
一锅肉多菜少的炖菜，多数晚上吃凉拌肉，
星期六吃煎羊肉……

\newsample

同上，但使用 |sameline|。

\samplesep

在拉曼恰的一个地方，其名我不愿提起，
不久前住着一位绅士
\begin{description}[font=\sffamily\bfseries, leftmargin=3cm,
    style=sameline]
\item[\textnormal{堂吉诃德最先有的}] 一支长矛,
\item[\textnormal{其次}] 一面旧盾,
\item[\textnormal{再次}] 一匹瘦马，
\item[\textnormal{最后，第四样}] 一只猎犬.
\end{description}
一锅肉多菜少的炖菜，多数晚上吃凉拌肉，
星期六吃煎羊肉……

\newsample

同上，但使用 |multiline|。注意如果条目正文太短。

\samplesep

在拉曼恰的一个地方，其名我不愿提起，
不久前住着一位绅士
\begin{description}[font=\sffamily\bfseries, leftmargin=3cm,
    style=multiline]
\item[\textnormal{堂吉诃德最先有的}] 一支长矛,
\item[\textnormal{其次}] 一面旧盾,
\item[\textnormal{再次}] 一匹瘦马，
\item[\textnormal{最后，第四样}] 一只猎犬.
\end{description}
一锅肉多菜少的炖菜，多数晚上吃凉拌肉，
星期六吃煎羊肉……

\fi

\normalsize

\section{后记}

\subsection{\LaTeX{} 列表}

众所周知，\LaTeX{} 预定义了三个列表：\texttt{enumerate}、\texttt{itemize} 和 \texttt{description}。这是一种非常常见的分类，在 HTML 中也能找到。然而，存在一个更通用的模型，基于三个字段——即标签、标题和正文——因此 enumerate 和 itemize 有标签（编号和不编号）但没有标题，而 description 有标题但没有标签。在这个模型中，你可以有一个带标签条目的 description，例如（当然，这个简单的方案远不能令人满意）：
\begin{verbatim}
\newcommand\litem[1]{\item{\bfseries #1,\enspace}}
\begin{itemize}[label=\textbullet]
\litem{堂吉诃德最先有的} 一支长矛，
……等等。
\end{verbatim}

\vskip6pt
\goodbreak
\hrule
\vskip6pt

\newcommand\litem[1]{\item{\bfseries #1,\enspace}}
在拉曼恰的一个地方，其名我不愿提起，
不久前住着一位绅士
\begin{itemize}[label=\textbullet]
\litem{堂吉诃德最先有的} 一支长矛,
\litem{其次} 一面旧盾,
\litem{再次} 一匹瘦马，
\litem{最后，第四样} 一只猎犬.
\end{itemize}

\vskip6pt
\goodbreak
\hrule
\vskip6pt

% This format in not infrequent at all and a tool for defining them is
% on the way and at a very advanced stage. It has not been include in
% version 3.0 because I'm not sure if the proper place is this package
% or \textsf{titlesec} and it is not stable enough yet.

\subsection{已知问题}

\begin{itemize}

\item
列表恢复基于环境名称，当 |\newenvironment| 包含列表时，你可能想使用 |\begin{<list>}| 和 |\end{<list>}|。使用相应的低级命令（即 |\<list>| 和 |\end<list>|）不是错误，但确保结果正确是你的责任。

\item boxed 标签（包括 enumerate 和 itemize）与 unboxed 标签混合使用的行为没有明确定义。boxed 和 unboxed 行内列表也是如此（甚至可能引发错误）。

\item 类似地，同时恢复系列和列表也是允许的，但同样其行为没有明确定义。

\item (3.5.2) 出现了与 2.x 的不兼容——如果你使用可选参数向 |\ref| 或其他需要可展开宏的宏传递值，会引发错误。快速修复是在 \texttt{before} 中让 |\makelabel| 指向 |\descriptionlabel|。

\end{itemize}

\subsection{3.0 新特性}

\begin{itemize}
\item 行内列表，带设置条目如何连接的键（即条目间的标点）。提供两种模式：|boxed| 和 |unboxed|。

\item |\setlist| 支持 \textsf{calc}（例如在循环中使用），你可以同时设置不同列表和层级。
\item 所有与标签相关的长度都可以取 |*| 值（不仅限于 |labelsep| 和 |leftmargin|）。其行为已保持一致，并有新值 |!|，它不计算最宽标签。

\item 使用 |\restartlist{<list-name>}|，可以重启列表计数器（在使用 |resume| 的情况下）。

\item |resume*| 可以与其他键组合使用。

\item 列表可以用系列全局收集，使它们被视为一个列表。要启动系列只需使用 |series=<series-name>|，然后用 |resume=<series-name>| 或 |resume*=<series-name>| 恢复。

\item "实验性的" |fullwidth| 已被新键 |wide| 取代。

\item|\SetLabelAlign| 定义新的对齐值。

\item 你可以定义"抽象"值（例如 |label=numeric|）和新键。
\end{itemize}

\begin{itemize}
\item (3.2) |start| 和 |widest*| 支持 \textsf{calc}。
\item (3.2) |\value| 可与 |widest*| 一起使用。
\item (3.2) 移除了 |\arabic| 等的一些内部限制。更灵活，代价是错误检查更"宽松"。
\end{itemize}
\subsection{错误修复}

\begin{itemize}
\item 星号值（例如 |leftmargin=*|）不能被覆盖，新值被忽略。
\item |nolistsep| 作为多个键中的第一个时不总是被识别，因此被当作短标签处理（即 |nol\roman*stsep|）。
\item |labelwidth| 不总是工作（当先前有 |widest| 和 |*| 时）。
\item 使用 |align=right| 时标签和后续文本可能重叠。
\item |description| 没有获得正确的列表层级。
\item 在某个版本（2.x?）|\value*| 停止工作。
\item (3.1) 不幸的是，\textsf{xkeyval} "杀死"了 \textsf{keyval}，因此后者已在 \textsf{enumitem} 中复制。
\item (3.3) 修复了一个严重错误——使用 |*| 时 |itemize| 和 |description| 都不工作。
\item (3.4) 修复了 boxed 模式中的间距问题（第一个条目前 |\unskip| 位置错误，条目间 spacefactor 错误）。
\item (3.4) |nolistsep| 未按预期工作，但由于这个错误已存在多年，提供了新键 |nosep|。
\item (3.4) |nolistsep| 与 |shortlabels| 的问题（见上文）并非在所有情况下都修复。希望现在修复了。
\item (3.5.0) 修复了与条目间 spacefactor 相关的问题。
\item (3.5.0) 修复了嵌套 boxed 行内列表的问题。
\item (3.5.1) \texttt{resume*} 只工作一次，后续的行为像 \texttt{resume}。
\item (3.5.2) 修复了不工作的 |\setlist*|。
\end{itemize}

\subsection{致谢}

我特别感谢 Lars Madsen 的评论和建议，他还发现了一些错误。

\end{document}

MIT License
-----------

Permission is hereby granted, free of charge, to any person obtaining a
copy of this software and associated documentation files (the
"Software"), to deal in the Software without restriction, including
without limitation the rights to use, copy, modify, merge, publish,
distribute, sublicense, and/or sell copies of the Software, and to
permit persons to whom the Software is furnished to do so, subject to
the following conditions:

The above copyright notice and this permission notice shall be included
in all copies or substantial portions of the Software.

THE SOFTWARE IS PROVIDED "AS IS", WITHOUT WARRANTY OF ANY KIND, EXPRESS
OR IMPLIED, INCLUDING BUT NOT LIMITED TO THE WARRANTIES OF
MERCHANTABILITY, FITNESS FOR A PARTICULAR PURPOSE AND NONINFRINGEMENT.
IN NO EVENT SHALL THE AUTHORS OR COPYRIGHT HOLDERS BE LIABLE FOR ANY
CLAIM, DAMAGES OR OTHER LIABILITY, WHETHER IN AN ACTION OF CONTRACT,
TORT OR OTHERWISE, ARISING FROM, OUT OF OR IN CONNECTION WITH THE
SOFTWARE OR THE USE OR OTHER DEALINGS IN THE SOFTWARE.
